\documentclass[12pt]{article}
\usepackage{graphicx} % Required for inserting images
\usepackage{authblk} % For affiliations
\usepackage{ebgaramond}
\usepackage{ragged2e}
\usepackage{hyperref}
\usepackage{amsmath}
\usepackage{amssymb}
\usepackage{makecell}
\usepackage{subcaption}
\usepackage{setspace}
\usepackage[a4paper, margin=1in]{geometry}
\usepackage{pdflscape}
\newcommand{\blandscape}{\begin{landscape}}
\newcommand{\elandscape}{\end{landscape}}

\hypersetup{colorlinks, anchorcolor = black, citecolor = blue, urlcolor=blue, filecolor=green, linkcolor = blue, plainpages = false}
\usepackage[skip = 5pt, indent=20pt]{parskip}

%TC:include \bibliography ref.bib
%TC:include \addbibresource{ref.bib}

\usepackage[
backend=biber,
style=apa,
sorting=nyt
]{biblatex}

\addbibresource{ref.bib}

\title{Education or Ethic? The Long-Term Impact of Protestantism on Electoral Behaviour in France}

\author[1]{Duarte Amaro}
\author[2]{Luis Sattelmayer}
\affil[1]{Department of Politics and International Relations, University of Oxford}
\affil[2]{Department of Political Science, University of Innsbruck}

\date{\centering Last updated: \today %\medskip \\ Word count: 13877 (10778 without references)
}

\begin{document}

\maketitle

\centering \large \textit{This is a draft. Please do not circulate.} \bigskip

\justifying \normalsize

\begin{abstract}
    Can a centuries-old religious divide still shape the political map of a secular republic? This paper examines the long-term impact of the Protestant Reformation on electoral behaviour in France. It revisits two competing legacies of Protestantism: an education mechanism, linking literacy and human capital to progressive political preferences, and an ethic mechanism, rooted in values of individual responsibility but also, in the French context, minority protection and constitutional liberalism. We use newly assembled, commune-level data on the Protestant population in France derived from an 1839 government-mandated survey of the non-Catholic churches. To our knowledge, this is the earliest available dataset on the Calvinist population in France at this level of granularity. This is combined with historical and contemporary data, including election results. Our findings show that the historical size of the Protestant population is negatively associated with support for the Radical Right, even controlling for education. This paper connects research on human capital and the Radical Right with recent work on how minority status can foster liberal commitments while also contributing to broader debates on the persistence of historical institutions and religious legacies.
\end{abstract}

\clearpage
\section{Introduction}

\setstretch{1.3}

\justifying

``Our Church, which knew once all the sufferings of persecution, feels ardent sympathy for your communities", wrote Marc Boegner, then President both of the Protestant Federation of France and of the National Council of the French Reformed Church, in a 1941 letter to the Chief Rabbi of France, following the adoption of the Vichy regime's anti-Semitic legislation (cited in \cite{leonardProtestantFrancais1955}, p. 267; see also \cite{cabanelMarcBoegnerOu2020}). 
More than eighty years later, ahead of the second round of the 2022 French presidential election, the Protestant Federation of France again intervened publicly, reminding its faithful ``how its commitments could be hindered in the event of a National Rally victory" \parencite{franceElectionsCommuniqueFPF2022}. These interventions are particularly succinct examples of a broader and longstanding association between Protestant milieus and comparatively liberal or anti-authoritarian political positions in France \parencite{latreilleForcesReligieusesVie1951, siegfriedGeographieElectoraleLArdeche1949}. 

Yet this association sits uneasily with dominant interpretations of Protestantism in the social sciences. Since Weber's classic thesis linking Protestantism to the ``spirit of capitalism'' \parencite{weberProtestantEthicSpirit2012}, Protestant culture has often been associated with individualism and weaker support for redistribution \parencite{bastenWorkEthicReligion2013, schaffUnequalSpiritProtestant2025}, which are, in turn, linked to support for right-wing parties. Moreover, much of the literature on National Socialism's electoral support and membership basis in Germany has emphasised the importance of Protestantism, whereas Catholic regions, especially where embedded in dense associational networks, proved more resistant to the radical right \parencite{brownNaziVoteNational1982,oloughlinElectoralGeographyWeimar2002,BRUSTEINW1995WJTN, haffertLongTermEffectsOppression2021, kolzer2025alte}. Protestantism therefore appears historically associated with both support for, and resistance to, authoritarian nationalism.

This article aims not only to revisit canonical works on geographical patterns of political behaviour and examine the association between historic Protestant settlement and support for the radical right in contemporary France, but also to explain this counter-intuitive relationship by reference to two competing mechanisms. We argue that the political consequences of Protestantism cannot be reduced either to theology or to a homogeneous set of cultural values. Instead, Protestantism's long-run effects might operate through human capital formation \parencite{beckerWasWeberWrong2009, mokyr2017culture, eire2016reformations}, which in turn foster depressed support for the radical right \parencite{abou-chadiBrahminLeftMerchant2021}, as well as through the historical experience of violent persecution \parencite{lupuLegacyPoliticalViolence2017}, with consequent preferences for minority protection and religious freedom \parencite{braunProtectorsPluralismReligious2019, rovnyEthnicMinoritiesPolitical2024}.

We leverage a novel dataset on the historic pattern of Protestant settlement in France at the commune level, derived from a government-mandated survey of non-Catholic churches conducted in 1839. The 1839 distribution of Calvinists predates modern party politics and mass electoral competition, and can therefore be treated as capturing the more original geographic footprint of the Reformation. In the post-revolutionary system of ``recognised faiths'', the French state financed, regulated, and monitored religious denominations, including the Reformed and Lutheran churches, resulting in the survey we employ here. To the best of our knowledge, this constitutes the earliest available dataset providing systematic information on the spatial distribution of the Calvinist population at such a fine level of geographical granularity. We combine this data with a range of historical and contemporary sources, including demographic indicators, educational outcomes, and electoral results.

The empirical analysis reveals a robust relationship between historical Protestant presence and contemporary electoral results. Communes with larger Protestant populations in the nineteenth century display somewhat higher levels of educational attainment today and significantly lower levels of support for the Radical Right. Importantly, the negative association between Protestant historical presence and Radical Right voting persists even after accounting for contemporary educational composition, suggesting that the relationship cannot be reduced entirely to differences in schooling levels alone.

The findings contribute to several strands of literature. First, they engage with research linking Protestantism to human capital formation by showing that the educational advantages associated with Protestant regions in the nineteenth century were modest by the 21\textsuperscript{st} century. Second, they speak to work on the social bases of Radical Right support by highlighting the enduring role of historical local contexts in shaping political attitudes and electoral behaviour. Finally, the article contributes to broader debates on the persistence of historical legacies. In the French case, the experience of religious minority status appears to have fostered not only distinctive educational trajectories but also durable commitments to political pluralism and resistance to exclusionary nationalism.

\section{Theory}

A growing literature in the social sciences treats historical experience as a central source of variation in contemporary political and economic outcomes \parencite{simpserDeadNotGone2018, abadHistoryNeverReally2021, charnyshHistoricalPoliticalEconomy2023}. 
%Past institutions, conflicts, and social hierarchies can persist through multiple channels, including institutional path dependence, local norms, social networks, collective memory, and intergenerational transmission \parencite{piersonPoliticsTimeHistory2004, wittenbergCruciblesPoliticalLoyalty2006, lupuLegacyPoliticalViolence2017}.
Canonical contributions link colonial institutions to state capacity and economic development \parencite{acemogluColonialOriginsComparative2001, nunnLongtermEffectsAfricas2008}, the slave trade to lower interpersonal trust \parencite{nunnSlaveTradeOrigins2011}, slavery to contemporary racial conservatism in the American South \parencite{acharyaPoliticalLegacyAmerican2016}, and anti-Semitic violence to persistent local patterns of exclusion \parencite{voigtlanderPersecutionPerpetuatedMedieval2012}. These legacies may even remain latent for long periods before re-surging in response to more permissive conditions, the increased salience of past events or of narratives that resonate with them, or the emergence of memory entrepreneurs who mobilise such narratives and connect present experiences to the local history -- especially when traditional alignments are upended \parencite{wittenbergCruciblesPoliticalLoyalty2006, dinasInterruptedContinuitiesLocal2024, ochsnerActivatedHistoryCase2024, foukaCollectiveRemembrancePrivate2023}.

% Studies of historical legacies include three components: an antecedent cause, a subsequent outcome, and a causal mechanism that connects the two \parencite{wittenbergConceptualizingHistoricalLegacies2015}. 
% Depending on whether the cause and the outcome are two instances of the same phenomena, \textcite{ruggeriLegaciesRepressionResistance2025} distinguish between legacy congruence and legacy alteration. The most straightforward type of legacy is persistence, when past behaviour, attitudes, or institutions are maintained until the present -- like lasting antisemitism in Germany, according to \textcite{voigtlanderPersecutionPerpetuatedMedieval2012}.
% %, or rightish voting in Hungary in Wittenberg's \parencite*{wittenbergCruciblesPoliticalLoyalty2006} account. 
% But these antecedent causes may also affect other kinds of outcomes: witnessing state violence against minorities not only shapes preferences for transitional justice but also lowers support for political parties associated with the perpetrators (\cite{charnyshGenocideProximityVictims2025, dejuanPoliticalEffectsWitnessing2025}, though see \cite{homolaLegaciesThirdReich2020}). 

%The mechanisms connecting cause and outcome can, in turn, be classified into institutions, coalitions, and cognitions \parencite{simpserDeadNotGone2018}. \textcite{brusteinSocialOriginsPolitical1988} distinguishes between structural and normative mechanisms.

While much of the work in the study of historical legacies has focused on the effect of regimes and formal institutions \parencite{popescuImperialBorderlandsInstitutions2023, doucettePreModernInstitutionsLater2024, guisoLONGTERMPERSISTENCELongTerm2016, kotkinHistoricalLegaciesCommunism2014} or conflict and violence \parencite{ruggeriLegaciesRepressionResistance2025, lupuLegacyPoliticalViolence2017, waldenHistoricalLegaciesPolitical, balcellsConsequencesVictimizationPolitical2012}, some attention has also been paid to the role of religion. Scholarship on the historical legacies of religious practice, and especially work in economics and economic history, has paid particular attention to human capital formation \parencite{beckerLutherGirlsReligious2008, beckerWasWeberWrong2009, cageLongTermEffectsPrinting2016, boppartPROTESTANTISMEDUCATIONREADING2014, gallegoChristianMissionariesEducation2010, hornungImmigrationDiffusionTechnology2014, squicciariniDevotionDevelopmentReligiosity2020, woodberryReligionSpreadHuman2011, mokyr2017culture, merton1936puritanism}. Other studies have focused on the effects of religion on local norms and individual preferences \parencite{bastenWorkEthicReligion2013, schaffUnequalSpiritProtestant2025, spaterProtestantEthicReexamined2019}, while yet another strand of the literature has emphasised the importance of religious organisations as key transmission mechanisms underlying persistence in political preferences \parencite{wittenbergCruciblesPoliticalLoyalty2006, haffertLongTermEffectsOppression2021}.

We connect this line of scholarship with (a) the political science literature on the increasing importance of education for party competition and (b) older works in French history and political sociology that emphasise the unique position of the French Protestant population as members not only as a religious denomination (overwhelmingly Calvinist) but also as a minority group of around 2 \%, as well as more recent studies that similarly emphasise minority groups' empathy towards each other. Two mechanisms should therefore connect the historical pattern of Protestant settlement in France with contemporary political outcomes, and especially with depressed support for the radical right: Protestantism should foster persistently greater human capital as well as stronger individual preferences for constitutional liberalism, freedom of religion, and minority protection.

\subsection{Education}

\subsubsection{Protestantism, Education, and Human Capital Formation}

A central argument in the literature is that Protestantism fostered human capital accumulation through its doctrinal emphasis on individual engagement with scripture \parencite{eire2016reformations, mokyr2017culture}. The principle of \textit{sola scriptura} encouraged literacy by making the reading of the Bible a religious obligation rather than a clerical monopoly\footnote{A point made in passing by \textcite[p. 151]{stinchcombeSocialStructureOrganizations1965}}. In this interpretation, Protestantism transformed incentives surrounding education and generated institutional support for mass schooling \parencite{straussTechniques1981, paulsen1885geschichte, stone1964educational}

\textcite{beckerWasWeberWrong2009} show that Protestant areas exhibited substantially higher literacy rates and economic development in 1871 Prussia, arguing that Protestantism exerted a causal effect on economic outcomes primarily through human capital accumulation. Protestant regions in Prussia already possessed denser educational infrastructure prior to industrialisation \parencite{beckerEffectProtestantismEducation2010} while Protestantism was also associated with narrower gender gaps in schooling and literacy \parencite{beckerLutherGirlsReligious2008}. Similarly, \textcite{boppartWhichConditionsDoes2013} show that Protestant exposure increased public expenditure on schooling and improved the educational performance of military conscripts in 19\textsuperscript{th} century Switzerland, with especially strong advantages in reading performance, emphasising the importance of scriptural literacy \parencite{boppartPROTESTANTISMEDUCATIONREADING2014}. 

Research on missionary activity in sub-Saharan Africa similarly finds that Protestant missions played a major role in the expansion of education and long-run human capital accumulation \parencite{woodberryMissionaryRootsLiberal2012, woodberryReligionSpreadHuman2011, cageLongTermEffectsPrinting2016} while the settlement of Huguenot refugees in Prussia generated substantial long-run productivity gains in textile manufacturing through the transfer of technical knowledge \parencite{hornungImmigrationDiffusionTechnology2014}. More broadly, the relationship between religious doctrine and educational investment is not unique to Protestantism. In 19\textsuperscript{th} century France, more intensely Catholic locations, measured through the share of the clergy that did not swear allegiance to the Civil Constitution of the clergy enacted by the revolutionary government in 1791, witnessed a slower adoption of the secular, technical curriculum  \parencite{squicciariniDevotionDevelopmentReligiosity2020}. Changes in Jewish religious institutions following the destruction of the Second Temple increased the importance of individual Torah reading and generated incentives for widespread education \parencite{botticiniJewishOccupationalSelection,botticiniFarmersMerchantsConversions2007}. Botticini and Eckstein argue that it was the exogenous change in religious dogma that prompted educational reforms and induced Jewish people to move into the cities, where the returns to their human capital were higher, in the eight and ninth century under the Abbasid caliphate. 

\subsubsection{Education in France}

The French case is perhaps best known as an illustration of the importance of central state action on educational expansion. Eugen Weber's \parencite*{weberPeasantsFrenchmenModernization1976a} classic account emphasised the state's role in extending schooling, raising literacy, and (perhaps most centrally) incorporating peripheral populations into a more homogeneous national political culture (see also \cite{furetReadingWritingLiteracy1983}). Public schooling and a state-designed curriculum helped standardise French as the national language, instil republican norms, and bind the local population into a shared national identity \parencite{GuillaumeblanccomFilesTheme}. However, educational expansion also depended heavily on local initiative, communal resources, and pre-existing social organisation. \textcite{furetReadingWritingLiteracy1983} show that school funding, teacher ability, and the curriculum all varied substantially, and often because of local demand for education. Weber's \parencite*[p. 303]{weberPeasantsFrenchmenModernization1976a} account similarly emphasises how the development of schooling was tied to the changing value and perceptions of its fruits. 

Geographically, the so-called Maggiolo line, stretching from Saint-Malo to Geneva, divided France into broad literacy zones up until the mid-19\textsuperscript{th} century: a more literate north and north-east, and a less literate west and south \parencite{furetReadingWritingLiteracy1983} -- precisely the area of greater Protestant settlement, as we will discuss below. By the late 18\textsuperscript{th} and early 19\textsuperscript{th} centuries, however, the south had began to catch up, and the old division was gradually replaced by a more pronounced east--west cleavage.

\subsubsection{The Persistence of Human Capital}

In turn, human capital exhibits substantial long-run persistence across generations. Using Swedish administrative data, \textcite{adermonDynasticHumanCapital2021} show that educational attainment persists strongly, especially if one considers the extended family, and that environmental factors explain at least 30 percent of the overall persistence. \textcite{blancChangePersistenceAge2020} find considerable intergenerational persistence in literacy in their study of Saint-Germain-d'Anxure between 1730 and 1895, even though mobility increased as literacy became more widespread. Regions formerly exposed to Jesuit missions before their expulsion from Latin America in 1767 continue to display significantly higher levels of educational attainment today \parencite{valenciacaicedoMissionHumanCapital2019} while Ottoman rule reduced long-run human capital accumulation with literacy differences persisting well into the twentieth century \parencite{popescuImperialRuleLongRun2022}. Similarly, \textcite{rochaHumanCapitalPersistence2017} demonstrate an increase in the literacy rate in municipalities that received immigrants with higher levels of schooling in early 20\textsuperscript{th} century Brazil; by the year 2000, individuals living in those municipalities still had 0.52 additional years of schooling, on average. The post-First World War partition of Togoland between the British and the French also led to persistent divergence, which \textcite{cogneauBordersThatDivide2014} attribute to the French administration's hostility towards missionary schools. 

At the same time, these historical legacies can fade \parencite{abadHistoryNeverReally2021}. \textcite{backhausFadingLegaciesHuman2019}, for example, show that large differences in school enrolment across the former partitions of Poland in 1911 gradually disappeared during the first two decades of Polish independence as national educational policies converged, so that literacy rates were uniform across each partition by the postwar period. Similarly, large educational disparities around the border of Portuguese-colonised Goa in 1961 vanished within thirty years \parencite{lehnerCultureInstitutionsFading}.

Taken together, this literature suggests that Protestantism encouraged literacy and educational investment, contributed to the development of local educational infrastructure, and fostered persistent human capital formation. Even where the original religious motivations weakened over time, the educational advantages they produced could remain, continuing to shape social and political outcomes.

\subsubsection{Contemporary negative association between education and support for the Radical Right}

In turn, vast scholarship sees education as a major pathway through which political preferences are built \parencite{abou-chadiBrahminLeftMerchant2021, gethin2022brahmin, kitschelt2023polarity, hooghe2025field, schafer2026education}. The rise of mass higher education in the second half of the twentieth century progressively restructured electorates around an educational divide and nowadays places education at the heart of political opinion formation \parencite{stubager2008education, stubager2009education, zollinger2026updating, l2026educational}. 

The literature finds a robust negative association between educational attainment and support for the Radical Right. Voters with lower education are overrepresented in Radical Right electorates across Western democracies, while voters with higher education concentrate in green, social-liberal, and new left parties \parencite{ivarsflaten2012voting, oesch2018electoral, abou-chadiBrahminLeftMerchant2021}. This divide has widened over time and now structures political competition as strongly as the classic class cleavage did in the mid-twentieth century \parencite{hooghe2026does}.

The mechanisms underlying this divide are multiple and reinforcing. Higher education appears to foster the values most directly opposed by the Radical Right -- tolerance of cultural difference, support for immigration, and cosmopolitan attitudes \parencite[see, for example,][]{cavaille2019education, stubager2008education}. Beyond level, the substantive \textit{field} of study matters too, as graduates of humanities, social sciences, and other adjacent disciplines are found to vote for GAL parties at substantially higher rates \parencite{hooghe2025field}. These effects are reinforced by educational self-selection earlier in life, by occupational sorting after graduation, and by the educational homogeneity of social networks. In this regard, higher- and lower-educated citizens increasingly form distinct social worlds, where educationally homogeneous networks reinforce political attitudes on both sides of the divide \parencite{de2025separated, schafer2026education}.

The overall finding of a negative association between higher education and Radical Right voting has direct implications for our argument regarding education and human capital. If Protestant historical presence fostered durable advantages in literacy and human capital that persisted into contemporary patterns of educational attainment, and if this educational attainment in turn robustly depresses Radical Right support, then we should see lower contemporary Radical Right vote shares in communes with larger historical Protestant populations. 

\subsection{Ethic}

\subsubsection{The Spirit of Capitalism}

Max Weber's \parencite*{weberProtestantEthicSpirit2012} seminal work on the Protestant ethic remains, over a century on, perhaps the first work to come to mind when considering the sociopolitical impact of the Protestant Reformation. While the work spurred by \textcite{beckerWasWeberWrong2009} has considered whether a ``Protestant ethic'' can instead be explained by differences in human capital accumulation, Weber's original argument was primarily cultural and psychological. He argued that certain forms of Protestantism, especially Calvinism, transformed attitudes towards work, discipline, and economic activity by recasting worldly labour as a religious vocation, a calling or \textit{Beruf}. The doctrine of predestination generated profound existential uncertainty, which believers sought to alleviate through signs of divine election manifested in Earth through disciplined conduct, sobriety, and economic success. Over time, this ethic encouraged systematic labour, frugality, and reinvestment rather than traditionalist consumption, contributing to the emergence of what Weber termed the ``spirit of capitalism'' \parencite{kalbergMaxWebersSociology2024}. Building on this tradition, \textcite{gorskiDisciplinaryRevolutionCalvinism2003} argues that Calvinism's most consequential legacy was a \textit{disciplinary revolution}. Through institutions such as consistories, communal surveillance, or reformed poor relief, Calvinist communities produced self-disciplined, orderly, and governable populations. This discipline, generated from below rather than imposed from above, could then be harnessed and rationalized by early modern states, lowering the costs of administration and contributing to state formation itself.



A colossal body of work across several disciplines has sought to empirically validate Weber's claims, reinterpret the mechanisms involved, or challenge the primacy he assigned to cultural values relative to institutions and education (for an overview, see \cite{marshallSearchSpiritCapitalism1993}). Recent work suggests the overall impact on economic development might be slight \parencite{delacroixBelovedMythProtestantism2001, cantoniEconomicEffectsProtestant2015a} while still finding support for an effect on cultural values and individual preferences. Exploiting the historical confessional border between the Swiss cantons of Vaud and Fribourg, \textcite{bastenWorkEthicReligion2013} show that Protestant Vaud displays systematically lower support for redistribution and government intervention. Protestant populations also industrialised earlier than neighbouring Catholic regions in Switzerland \parencite{spaterProtestantEthicReexamined2019} while the Reformation exacerbated inequality in early modern Germany through the introduction of more particularistic forms of poor relief \parencite{schaffUnequalSpiritProtestant2025}. 

\subsubsection{French Protestants and Minority Status -- a \textit{liberal ethic}?}

Much of the quantitative literature on the consequences of the Reformation has focused either on cross-national differences or on contexts in which Protestants formed a majority, or at least a politically substantial minority. France presents a different case. French Protestants were not simply adherents of a religious doctrine associated with literacy, discipline, or economic individualism; they were also a small minority which underwent centuries of violent persecution. We argue that the historical experience of being a vulnerable minority in a predominantly Catholic political order has led them to develop political orientations shaped by minority status, chiefly scepticism towards homogenising claims of national unity and support for institutions that protect pluralism and freedom of religion, both of which should impact political behaviour today, more so than a Weberian work ethic.

We build on Rovny's \parencite*{rovnyEthnicMinoritiesPolitical2024, rovnyAntidoteBackslidingEthnic2023} argument that minority groups can become ``circumstantial liberals” who support liberal institutions not because of doctrinal or ideological reasons but because those institutions constrain majority rule. Communities that cannot expect to control the state, or to leave it altogether for an ethnic kin one, have strong incentives to defend civil liberties, minority rights, constitutional limits, and political pluralism, and should thus be less likely to support parties which claim to speak for a culturally homogeneous national majority. 

In a similar vein, \textcite{braunReligiousMinoritiesResistance2016, braunProtectorsPluralismReligious2019} argues that religious minorities are more likely to protect persecuted groups during episodes of mass killings or genocide. Minority communities had not only good reason to recognise the danger of exclusionary rule and emphathise with the targets of persecution but also the organisational capacity (through their churches, clergy, associations, and trust networks) to coordinate assistance. Using geocoded data on Jewish evasion in Belgium and the Netherlands during the Holocaust, Braun shows that proximity to Catholic churches increased evasion in predominantly Protestant areas whereas proximity to Protestant churches had the same effect in predominantly Catholic areas. Assistance to threatened neighbours was the result not of different confessional doctrines or inherent features of each religious group, but of their local position and minority status.

Minority liberalism can persist even once the original conflict has faded, as long as organisational strength and community bonds ensure their social transmission \parencite{charnyshRoleCommunitiesTransmission2022, cremaschiPoliticalLegaciesWartime2025}. As \textcite{haffertLongTermEffectsOppression2021} shows, nineteenth-century Prussian repression of Catholicism in Germany stimulated greater mobilisation and spurred the creation of lay organisations. Catholic regions where oppression was more intense and where the Catholic milieu was better developed consequently display lower levels of support for the radical right \textit{Alternative für Deutschland} today (though see \cite{arzheimerOppressionCatholicsPrussia2024, manowErbeKulturkampfsOder2023}). Similarly, \textcite{wittenbergCruciblesPoliticalLoyalty2006} highlights the importance of local church-based social networks in explaining the persistence of right-wing support in Hungary before and after communism. Churches mattered because they embedded individuals in local networks, fostered within-group cohesion, and transmitted political loyalties, insulating churchgoers from the surrounding socialist milieu. In a certain sense, the disciplinary institutions that Gorski \parencite*{gorskiDisciplinaryRevolutionCalvinism2003} identifies as instruments of state rationalization where Calvinism was dominant played a different role in France. Rather than serving state power, consistorial discipline and dense communal organisation sustained a persecuted minority through clandestinity, providing precisely the organisational capacity that later work identifies as central to the transmission of minority political orientations \parencite{braunProtectorsPluralismReligious2019, charnyshRoleCommunitiesTransmission2022}.

%\textcolor{red}{Collective memories also have to be mobilised \parencite{cremaschiPoliticalLegaciesWartime2025, foukaCollectiveRemembrancePrivate2023, dinasFamilyHistoryAttitudes2021}}

There is substantial historical evidence that French Protestantism's historical persecution and minority status persisted in its collective memory long after their emancipation. Perhaps most famously, Siegfried (\cite*{siegfriedGeographieElectoraleLArdeche1949}; Siegfried in \cite{latreilleForcesReligieusesVie1951}) argued that the persecutions of the seventeenth and eighteenth centuries left French Protestants with a vivid memory of their vulnerability. This memory generated a deep mistrust of authoritarian or anti-liberal regimes, an instinctive defence of minority rights, and a strong commitment to Republican principles as a guarantee of religious liberty. French Protestant republicanism was not reducible to social class or economic interest, and indeed often went against it. The Revocation of the Edict of Nantes and the experience of the \textit{Désert} \footnote{The period after 1685 when French Protestants, deprived of legal recognition, were forced into clandestine worship.} anchored Protestants on the left for as long as monarchy and the right remained associated with Catholic power and hostility to a minority-protecting \textit{laïcité} \parencite{cabanelGirondeGardoiseExistetil2004, encreveProtestantsFrance18001985, schramProtestantismPoliticsFrance1954}.

Indeed, Calvinist opposition to religious uniformity arguably formed part of a longer genealogy of Republicanism in France. \textcite{vankleyReligiousOriginsFrench1996} argues that the French Revolution was shaped not only by secular Enlightenment ideas but also by earlier religious conflicts that weakened sacral monarchy and furnished languages of resistance to absolutism. By the time of the Third Republic, religion was among the key private determinants of republican affiliation among several founding republican intellectuals, whose social worlds were often suffused with Protestantism \parencite{hazareesinghIntellectualFoundersRepublic2004, cabanelGirondeGardoiseExistetil2004}. Protestant networks helped supply part of the Republic's intellectual and administrative infrastructure, but this visibility also made Protestants targets of anti-republican suspicion. In Catholic nationalist and radical right extremist thought, associated with Charles Maurras, the Protestant appeared alongside the Jew, the Freemason, and the Socialist as one of the four symbolic internal enemies of the French nation \parencite[87, 134]{hazareesinghPoliticalTraditionsModern1996, schramProtestantismPoliticsFrance1954}. Echoing Braun's \parencite*{braunProtectorsPluralismReligious2019, braunReligiousMinoritiesResistance2016} findings, Protestants made up only 2\% of the French population in 1940 but represented over 10\% of all French who rescued Jews during the Holocaust. The experience of persecution had fostered strong affinities between the two religious minorities \parencite{cabanelProtestantMinoritiesJudeoProtestant2014}. 

This literature leads us to expect that, because French Protestants were historically a persecuted minority, their communities should have developed a durable attachment to constitutional liberalism, minority rights, and freedom of religion. These orientations, if persistent, should relate to lower contemporary support for the Radical Right. We therefore expect historical Protestant presence to be associated not only with higher human capital, but also with lower contemporary support for the Radical Right, even net of education, insofar as Protestant minority status fostered a lasting liberal ethic.

\subsection{Alternative Accounts}

Although we consider these two theoretical mechanisms compelling, we discuss several alternative explanations here. If the literature has devoted increasing attention to the study of historical legacies, it has also vigorously discussed their theoretical cogency and empirical robustness \parencite{abadHistoryNeverReally2021} and drawn scholars' attention to more temporally proximate factors as well as to regional specificities \parencite{arzheimerOppressionCatholicsPrussia2024, manowErbeKulturkampfsOder2023, edwardsDidProtestantismPromote2021,pepinskyModelingSpatialHeterogeneity2024}

\subsubsection{Protestant Settlement Patterns and Self-Selection}

A possible alternative explanation is that Protestantism simply took root in places that were already more educated, commercially developed, or predisposed towards political dissent. While the geography of French Protestantism was not random, the historical evidence suggests that the durable settlement pattern of French Protestantism was shaped less by pre-existing levels of human capital than by political and military factors. In the Holy Roman Empire, the adoption of Protestantism could reflect the choices of local rulers, under \textit{cuius regio, eius religio}, and therefore pre-existing characteristics of their territories \parencite{cantoniAdoptingNewReligion2012a}. \textcite{grzymala2024popes} shows that the Reformation spread and persisted not where local conditions were initially most favourable to it, but where the fragmentation of political authority allowed sympathetic local rulers to shelter it from both papal and imperial retribution and repression. In France, by contrast, the monarchy remained Catholic and hostile to religious pluralism at the same time that it consolidated its hold over the national territory. The geography of French Protestantism therefore reflects the uneven reach of royal and judicial repression more than early diffusion, which was not, at any rate, radial or linear \parencite[p. 11]{leonardProtestantFrancais1955}. Where Protestantism in the Holy Roman Empire survived \textit{because} no central authority could suppress it, the French case shows Protestantism surviving only where a centralising but still uneven Catholic state had not yet succeeded in suffocating it at the root.

The initial spread of the Reformation was considerably broader than the later map of French Protestantism would suggest. \textcite{bostCarteProtestantismeFrancais1924} emphasises that Lutheran and Reformed ideas initially circulated across much of the kingdom, especially along major communication routes. Before 1541, denunciations of heresy appeared not only in the South but also in central and northern cities such as Meaux, Caen, Amiens, Orléans, Bourges, and Tours. The Wars of Religion completely transformed the geographical distribution of French Protestantism. Protestant communities in regions closer to Paris and its \textit{Parlement} were more exposed to repression and military defeat (\cite[p. 17]{leonardProtestantFrancais1955}, \cite{benedictFaithFortunesFrances2001, moursEglisesReformeesFrance1958}). In \textcite{mannAutonomousPowerState1984}'s terms, the geography of early Protestant survival reflected the uneven infrastructural power of the early modern French state: Protestantism endured more readily where royal coercion was weaker, delayed, or militarily contested. It is worth noting that the areas of surviving Protestant settlement were (with the exception of Alsace) south of the Maggiolo line, where human capital had historically been lower, further dispelling concerns about self-selection.

The Edict of Nantes in 1598 locked in this geography. Although it granted a limited measure of toleration to Huguenots, it did not create a general right to Protestant worship throughout the kingdom. Rather, it recognised and regulated worship in places where Protestantism had already secured a presence, while Catholicism remained the established religion of the kingdom \parencite{bostCarteProtestantismeFrancais1924, moursEglisesReformeesFrance1958}. The settlement therefore institutionalised the results of prior conflict: areas where Protestant communities had been destroyed or weakened before 1598 were unlikely to recover, while areas where they retained military or organisational strength were able to preserve recognised places of worship \parencite[p. 119]{benedictFaithFortunesFrances2001}. Subsequent repression under Louis XIII and Louis XIV, culminating in the Revocation of the Edict of Nantes in 1685, reinforced this pattern. The Revocation abolished the legal protections previously afforded to Protestants and produced widespread persecution and emigration, despite formal prohibitions on leaving the kingdom. \textcite{bostCarteProtestantismeFrancais1924} argues that this process again deepened the North--South divide. In regions where Protestant communities were already small or isolated, emigration and conversion often led to their disappearance \parencite{benedictFaithFortunesFrances2001, benedictPopulationReformeeFrancaise1987}. By contrast, areas with denser Protestant communities were better able to preserve religious identity through the clandestine organisation.

By the late modern and contemporary periods, the geographical distribution of Protestantism still largely reflected the outcome of the Wars of Religion rather than sixteenth-century educational advantage. Nineteenth century migration did not largely modify this, despite Flemish, American, and, after 1870, Alsatian Protestant settlement in cities such as Bourges, Reims, and Nancy \parencite{bostCarteProtestantismeFrancais1924}. 

\subsubsection{Structural Conditions and Modes of Production}

Even if not directly shaping Protestant settlement, one might nevertheless argue that, given the concentration of the Protestant population in the South, any political outcome might be due more to regional specificities, particularly as regards economic structure, than to Protestantism. Indeed, one of most important accounts of French electoral geography centres on the type of soil, population density, and the mode of land tenure, which combined to result in different degrees of peasant dependence on the church and the nobility and, in turn, in greater or lesser support for the Left \parencite{siegfriedTableauPolitique1964}. Focusing on Mediterranean France, \textcite{loubereRadicalismMediterraneanFrance1974} emphasises the local economic structure in explaining patterns of political support. Building on these, \textcite{brusteinRegionalModeofProductionAnalysis1981a, brusteinSocialOriginsPolitical1988} explains persistent regional voting patterns through distinct modes of production: a commercially oriented north-east based on cash tenancy and wage labour, a subsistence western zone based on share-cropping and tenancy, and a market-oriented Mediterranean zone based on small-owner cultivation or wage labour. These regional socioeconomic structures correspond to enduring political alignments, with the west tending towards the Right and the Mediterranean and north-east more often towards the Left. 

Yet all of these scholars also admit the important role of religion. Siegfried does so extensively in other work (e.g., \cite{siegfriedGeographieElectoraleLArdeche1949}) while \textcite{loubereRadicalismMediterraneanFrance1974} considers the numerous exceptions to the general trend towards increased radicalism prompted by the successive shocks that affected Southern winemakers, chief among them local religious composition. Brustein is a strong opponent of accounts suggesting an independent role for individuals' shared norms, values, or beliefs, which he considers not only ``impossible to refute" but also appearing to ``strain logic in supposing that individuals’ votes reflect nonrational factors such as tradition rather than self-interest" (\cite{brusteinSocialOriginsPolitical1988}, p. 18). He nevertheless finds that both the mode of production \textit{and} religiosity are significant predictors of voting for the Left in 1981, and that the two variables together explain more of the variation that either does separately. Not only do these structural accounts make some concessions to the role of religion, but they are also generally captured by the inclusion of department fixed effects across most of our model specifications.

\subsubsection{Religious Competition and Education Provision}

An alternative explanation for the relationship between Protestant settlement and current educational outcomes could be that increased religious competition, and the wider dogmatic shift following the Counter-Reformation, spurred the Catholic Church to develop its educational provision.  Indeed, \textcite{furetReadingWritingLiteracy1983, furetThreeCenturies} argue precisely that, under the \textit{Ancien Régime}, it was the Church, and not the State, which took up education as a cause. In their words, the Church "was forced to fight the Reformation with the weapons of the Reformation" \parencite[p. 60]{furetReadingWritingLiteracy1983}, although the focus remained on reading rather than writing. It was the French Revolution which brought about a conflict between the two anchored on education, and in which their roles were reversed. 

However, while that was the overall stance of the pre-Revolutionary Catholic Church, the actual provision of education was much more dependent on local conditions. Schools (and schoolteachers), when they existed, were often financed directly by the families, the wider community, or through village revenues. Furet and Ozouf provide extensive evidence on the idiosyncrasies and unevenness of funding, teacher ability, curricula, and the school calendar -- and that is is this uneven demand which best accounts for the uneven spread of French schooling on the eve of the Revolution (p. 82; see also \cite{weberPeasantsFrenchmenModernization1976a}, pp. 304-307). At any rate, if our results are owed to religious competition and increased provision of \textit{Catholic} schooling in Protestant areas, then the effect of Protestant settlement on the contemporary share of people with higher education should be moderated by the Catholic religiosity (measured through the share of priests pledging allegiance to the Civil Constitution in 1791).

\subsubsection{Linguistic Minority Status}

A further alternative explanation is linguistic minority status. Centre-periphery conflict emerges where culturally or linguistically distinct regions resist national integration \parencite{lipsetCleavageStructuresParty1967}. State-led schooling and nation-building were central to the replacement of regional languages by French \parencite{weberPeasantsFrenchmenModernization1976a}, and Protestant settlement could have been concentrated in areas where French was historically weaker. Recent work by \textcite{ziblattWealthTonguesWhy2023} shows that historically peripheral linguistic regions are more likely to support the radical right in Germany, which would mean that our estimate of a negative association between Protestantism and radical right support in France would actually be biased downward. In any event, we control for local linguistic minority status, measured as the share of French speakers in 1886.

\subsubsection{Protestant Voting or Christian Voting?}

A further alternative explanation is that communes with historically strong Protestant settlement are also more religious today, and that religiosity itself (rather than denomination) reduces support for the radical right. \textcite{arzheimerChristianReligiosityVoting2009} find that Christian religiosity has few direct effects on radical-right voting in Western Europe; religious voters are not necessarily more xenophobic, but are often retained by Christian Democratic parties. \textcite{marcinkiewiczReligiousVotersSupport2022} likewise show that the association between religiosity and populist radical right support varies by region, with opposite effects in Western and Central and Eastern Europe. In France, while Catholic religiosity, measured through church attendance, is still strongly related to support for the mainstream right, it is insignificantly or even negatively related to support for the radical right \parencite{lewis-beck2004, boyLelecteurSesRaisons1997a, mayerFrontNationalVote1998, michelatSystemeSymboliqueCatholique2015, overosDieuDansCommune2023}. We attempt to disentangle these two mechanisms by examining also the relationship between Protestant settlement and support for the mainstream Right, which we expect to be weaker than that for the Radical Right while still negative, given the mainstream Right's connotation with the Catholic Church.

% Individual-level data?

\section{Data and Methods}
The empirical core of this paper is a novel dataset on the historic Protestant population at the communal level. We hand-transcribed it from a digitalized government-mandated survey of the non-Catholic recognized churches which was run in 1839. To the best of our knowledge, this is the earliest available source which provides systematic information on the spatial distribution of French Protestants at the level of the commune. 

Our choice of historical data is driven by both theoretical considerations and practical constraints. Our argument concerns the long-run imprint of the Reformation on human capital formation as well as the specific experience of religious minority status. What matters for our purposes is where Protestant communities originally took root, survived persecution, and built social networks over centuries, not where Protestants happen to live today.\footnote{Even though we would expect these two measures to be correlated.} Contemporary Protestant population figures would reflect nearly two centuries of demographic and sociological change and would be difficult to interpret as a trace of the Reformation's original geographic footprint. On the practical side, systematic commune-level data on the contemporary Protestant population in France does not, to our knowledge, exist. Unlike the nineteenth-century system of \textit{cultes reconnus}, which generated detailed administrative records as a byproduct of state regulation of recognised faiths, no equivalent source tracks the spatial distribution of French Protestants at this level of granularity today. In fact, French law prohibits surveying about religious denomination ("Loi informatique et liberté", loi no 78-17 du 6 janvier 1978 relative à l'informatique, aux fichiers et aux libertés).\footnote{https://www.legifrance.gouv.fr/loda/id/LEGITEXT000006068624} The 1839 survey thus provides the closest available approximation to the geographic distribution of Calvinist settlement that emerged from the Wars of Religion and consolidated under the \textit{Ancien Régime}.

In itself, the survey is the product of the post-revolutionary system of \textit{cultes reconnus}, first instaured by the Concordat of 1801 and the Organic Articles of 1802, under which the French state regulated and monitored the activities of officially recognised religious denominations, which from 1802 onwards included the Reformed and Lutheran churches alongside Roman Catholicism (Judaism was added in 1808). Pastors of the recognised Protestant churches were salaried public officials, and their parishes were subject to administrative oversight by the Ministry of \textit{Cultes}. Under the Restoration and the July Monarchy, the systematic collection of information on Protestant communities formed part of a broader effort to govern religious diversity, allocate state resources between recognised faiths, and integrate historically distinct confessional populations into a centralised national administration. The 1839 survey is the most detailed of these enumerations and was prepared in close cooperation with the consistorial authorities, who provided the underlying figures on parish boundaries, population, and worship.

The unit of observation in the original survey is the parish (\textit{section}), which is in turn nested within a local consistory (\textit{consistoire}, the rough Protestant equivalent of a diocese\footnote{It is worth noting that the consistory was introduced by the State through the Organic Articles and did not correspond to previous forms of organisation, which instead centred on a pyramid of assemblies (each local church's council of elders would elect delegates to a colloquium representing several local churches, which in turn would send representatives to a provincial synod, whose representatives would then make up the general synod). \textcite{encreveProtestantsFrance18001985} claims that the consistorial system created a hierarchy between local churches which was anathema to the original organisation of French protestantism, though \textcite{schramProtestantismPoliticsFrance1954} argues that the new system matched the wishes of the more bourgeois Protestants.}) and a main church (\textit{église consistoriale}). Each parish entry lists the communes that fall within its territorial jurisdiction, together with, for each commune: the type of place of worship (\textit{temple}, private house, public building, open air, etc.); the resident population; the resident Protestant population; the total Protestant population of the parish; the distance between the commune's center (\textit{chef-lieu}) and that of the parish; the place of residence of the pastor; and the name of the pastor. Figure \ref{fig:example} in the Appendix shows one of the survey's sheets as an example. The documents transcribed so far cover only the Reformed (Calvinist) churches; the equivalent enumeration of the Lutheran communities, concentrated in Alsace and the Montbéliard region, will be incorporated in subsequent work and is, at the time of writing, not included in the empirical analysis.

The 1839 document was transcribed by hand into a spreadsheet preserving the original parish-commune structure, yielding 2{,}455 row-level observations across 96 consistories and 112 main churches, covering all metropolitan departments. Each commune was matched to its present-day INSEE\footnote{The French national bureau of statistics.} identifier using the Cassini database of historical French communes \parencite{LecfcfrNewArticles}\footnote{Available at \url{http://cassini.ehess.fr/fr/html/6_index.htm}.}. The resulting dataset provides information on the Protestant population of 2{,}408 distinct communes for which the underlying source records a determinate value.

Several types of irregularities in the original document required additional cleaning before the data could be used for statistical analysis. First, in a small number of cases, the survey reports an aggregate number of Protestants spread across an unspecified set of ``surrounding communes'' or ``neighbouring hamlets'' without identifying them individually. Where the geographical referent could be reasonably reconstructed from the context (e.g., the surrounding communes of Tournon-sur-Rhône in Ardèche, or the communes of the canton of Rocroi in the Ardennes), we distributed the aggregate Protestant population across the relevant communes proportionally to their 1839 resident population, drawing on the historical population estimates of \textcite{cageHistoryPoliticalConflict2025}. Where no such reconstruction was feasible (e.g., for the Parisian metropolitan area, where the document explicitly acknowledges the difficulty of enumeration; for the suburbs of Bordeaux; or for unidentified communes in the departments of Ain, Aisne, Calvados, Manche, Doubs, Loiret, and Dordogne), the affected communes were assigned missing values rather than zero, to avoid spurious extrapolation. A full overview of the irregularities and a detailed list of our decisions can be found in Section \ref{sec:dataset} of the Supplementary Material. 

Second, despite an overwhelming and impressive coherence and consistency in French communal organisation since the 19\textsuperscript{th} century, a number of commune-level entries refer to administrative units that have since been merged, dissolved, or reorganised otherwise. These cases were harmonised to the current set of communes and the Protestant population was summed across the constituent units of the present-day commune. Third, where the same commune appears in more than one row, we adopted a conservative approach. When the two entries reported the same population and similar Protestant counts, we treated the case as double-counting and retained the higher of the two figures in areas where the Protestant population was lower and the lower figure in areas of greater Protestant settlement. When the entries clearly refer to distinct sections of the commune, we added them up into one observation. Finally, and arguably most importantly, for communes that were not explicitly mentioned in the original 1839 survey but were located in the departments that it covers, we imputed a zero value for the Protestant population.\footnote{An exception being the departments listed above, for which the survey itself signals incomplete coverage.} We use a marker to distinguish communes whose values were taken directly from the original source (“Original”) from those that were interpolated under the assumptions described above (“Extrapolated”). In our empirical analysis, we conduct robustness tests for both subsets. Figure \ref{fig:map_own} shows the Calvinist share of the population in 1839 per commune. For comparison, Figure \ref{fig:fischbacher} shows a 1923 map of the Protestant population in France. There is a strong overall correspondence in Protestant hotspots, even if the 1923 map shows a smoother and more diffuse distribution.

\begin{figure}[h]

\begin{subfigure}[c]{0.55\textwidth}
\includegraphics[width=0.9\linewidth]{map.jpeg} 
\caption{\Centering Distribution of the Calvinist population per commune in 1839}
\label{fig:map_own}
\end{subfigure}
\begin{subfigure}[c]{0.45\textwidth}
\includegraphics[width=0.9\linewidth]{carte.jpg}
\caption{\centering Distribution of the Protestant population in 1923 \parencite{louisCarteFranceProtestante1923}}
\label{fig:fischbacher}
\end{subfigure}

\caption{Distribution of the French Protestant Population}
\label{fig:maps}
\end{figure}

\subsection{Additional Historical and Contemporary Controls}

The 1839 communal Protestant share is combined with a range of historical and contemporary commune-level sources. The electoral outcomes for the legislative elections of 1993, 1997, 2002, 2007, 2012, 2017, and 2022 are drawn from \textcite{cageHistoryPoliticalConflict2025} and form the dependent variables of our first models, with attention to the first-round vote shares of the Radical Right (the \textit{Front National}, renamed \textit{Rassemblement National} in 2018) and of the mainstream right (the \textit{Rassemblement pour la République}, succeeded by the \textit{Union pour un mouvement populaire} and \textit{Les Républicains}). 

Furthermore, we include a vector of contemporary control variables, all coming from \textcite{cageHistoryPoliticalConflict2025}. We add controls for today's levels of communal mean income, the total resident population, the share of non-French citizens, the unemployment rate, mean property prices, and the share of high-school diploma (\textit{baccalauréat}) and tertiary education holders.

Next, we introduce a number of historical controls to address the principal alternative explanations discussed above. We capture pre-existing levels of human capital by the share of married individuals signing their marriage contract in 1686 as a proxy of literacy. \textcite[pp. 9-18]{furetReadingWritingLiteracy1983} validate this measure by comparing it to self-declared literacy in the national census from 1866 onwards and to the self-declared literacy of male military conscripts (itself verified by the Ministry of War in an 1864 examination of conscripts' literacy), finding that ability to sign one's name at marriage indeed captures the capacity to read and write. Pre-revolutionary religiosity is captured by the share of constitutional priests in 1791, while early levels of economic development are operationalized through the ratio of commune to national revenue in 1790. All three variables are again drawn from \textcite{cageHistoryPoliticalConflict2025}. 

Additionally, we build on \textcite{albertusStateBuildingRebellionRunUp2025} and their study on state building and rebellious efforts in revolutionary France. We include indicators of pre-1789 state capacity -- distance to the nearest Cassini road, distance to the nearest postal relay station in 1708, the presence of a letter post in 1710, and the presence of a police (\textit{maréchaussée}) brigade in 1720. Local exposure to early modern collective violence is measured by the count of rebellions, riots, and other collective events recorded in the Jean Nicolas database between 1661 and 1789 \parencite{gayvictorJeanNicolasDatabase2025}, with separate counts by intensity level and by whether the original record identifies a Protestant presence. We use the share of French speakers in 1886, sourced from the Historical Ethnic Geography data \parencite{muller-creponRightPeoplingStateNationalism2025}, as a measure of local linguistic diversity. Finally, we also add geographical controls, mean altitude and the geographical coordinates of the commune's centroid, which are drawn from publicly available French administrative data.

\subsection{Empirical Strategy}

\subsubsection{Baseline OLS}

Our baseline specification regresses commune-level outcomes -- the first-round vote share of the Radical Right and, in mediation analyses (see below), the share of the population with a higher education degree -- on the 1839 Protestant share, pooling the seven legislative elections between 1993 and 2022 with robust standard errors clustered at the commune level:

\[
y_{it} = \beta\,\text{Protestants1839}_{i} + \mathbf{X}_{it}'\boldsymbol{\gamma} + \mathbf{Z}_{i}'\boldsymbol{\delta} + \alpha_{g(i)} + \tau_t + \varepsilon_{it},
\]

where $\mathbf{X}_{it}$ is a vector of contemporary commune-level controls, $\mathbf{Z}_i$ is a vector of historical controls, $\alpha_{g(i)}$ are geographic fixed effects, and $\tau_t$ are election-year fixed effects.

We build up the specification in four steps. We begin with department fixed effects, which restrict identification to within-department variation across communes and absorb any time-invariant feature of regional political culture, religious tradition, or institutional history that might vary at the departmental level. To the department fixed-effects specification we sequentially add the vector of historical controls -- pre-existing literacy, revolutionary religiosity, early levels of economic development, pre-1789 state capacity, early modern collective violence, linguistic diversity, and geography -- and then the vector of contemporary controls -- income, population, the share of foreigners, unemployment, property prices, welfare receipt, and the socio-professional composition of the labour force. The progressive saturation of the specification allows us to assess whether the relationship between historical Protestant presence and contemporary Radical Right voting survives increasingly demanding adjustments for both deep-rooted and present-day differences between communes.

However, contemporary education is a post-treatment variable: if Protestant settlement affected local educational trajectories over the long run, then conditioning on current education blocks part of the very effect we seek to explain. More importantly, education is likely to be jointly determined with other contemporary characteristics such as income, unemployment, population size, and housing markets, many of which may themselves be downstream consequences of historical Protestant presence. Controlling for these variables may induce post-treatment bias. We therefore complement the baseline regressions with sequential g-estimation, following \textcite{acharyaExplainingCausalFindings2016}. Sequential g-estimation is designed precisely for settings in which a treatment may affect both a mediator and the confounders of the mediator--outcome relationship. 

We take latitude, longitude, mean altitude, and the department as baseline covariates and implement the procedure in four variants. First, we treat the historical controls as intermediate variables and estimate whether Protestant share predicts Radical Right voting once the variation associated with these historical pathways is held fixed\footnote{Although these variables are measured before the 1839 survey from which our measure of the Protestant population is derived, we are wary of regarding them as pre-treatment. Historical evidence suggests that the broad geography of French Protestant settlement was largely established during the Wars of Religion and consolidated by the Edict of Nantes in 1598. Variables such as marriage literacy, revolutionary religiosity, local wealth, and patterns of collective contention may therefore themselves reflect the long-run consequences of Protestant persistence rather than exogenous differences that preceded it. We consequently treat these variables as potential intermediate variables and interpret specifications that condition on them as estimating effects net of important historical pathways rather than total effects of Protestant settlement.}. Second, we repeat the exercise using only the contemporary controls as intermediate variables, asking whether the association persists net of present-day socioeconomic conditions. Third, we include both historical and contemporary controls as intermediate variables, yielding a particularly demanding estimate of the residual association that remains after accounting for both sets of post-treatment pathways. In these three specifications no mediator is specified in the blip-down step. Finally, we add contemporary higher education as the blip-down variable while treating the historical controls, the contemporary controls, or both as intermediate confounders. In this case, sequential g-estimation removes the estimated contribution of higher education to the outcome and estimates the average controlled direct effect of Protestant share when higher education is fixed at its mean level \parencite{acharyaExplainingCausalFindings2016}. We implement these estimators using the \texttt{sequential\_g()} procedure provided by the \texttt{DirectEffects} package, which operationalises the sequential g-estimation framework of \textcite{vansteelandtEstimatingDirectEffects2009} and \textcite{acharyaExplainingCausalFindings2016} for the estimation of average controlled direct effects.

\subsubsection{Mediation Analysis and Robustness Checks}
Our theoretical framework identifies two complementary pathways through which historical Protestant presence may shape contemporary support for the Radical Right: the accumulation of human capital and an ethic rooted in the historical experience of religious minority status and its associated commitments to constitutional liberalism and minority protection. The sequential g-estimation analyses described above allow us to assess whether the association between Protestant share and Radical Right voting persists once the education channel is held fixed. However, they do not provide a decomposition of the total association into educational and non-educational components. To obtain an indicative measure of the relative importance of these pathways, we therefore conduct a supplementary causal mediation analysis following \textcite{imai2010general}, implemented in R using the \texttt{mediation} package \parencite{tingley2014mediation}.

In our models, the mediator is the contemporary share of the resident population holding a higher education degree. We estimate two equations: 1) a mediator model regressing the share of higher education holders on the 1839 Protestant share and the full vector of controls, and 2) an outcome model regressing the vote share for the Radical Right on the 1839 Protestant share, the mediator of education, and the same controls. From these, we dissect the total effect of the 1839 Protestant share into an average causal mediation effect (ACME), which captures the portion of the effect that operates through contemporary accumulated human capital, i.e. education, and an average direct effect (ADE), which captures the part of the variance that operates through all other channels. These other possible channels include, in our interpretation, the persistent influence of minority-derived political orientations following the argument by \textcite{rovnyAntidoteBackslidingEthnic2023, rovnyEthnicMinoritiesPolitical2024}. 

Unlike the controlled direct effects estimated via sequential g-estimation, however, this decomposition relies on the assumption of sequential ignorability \parencite{imai2010general, imai2010identification}, which here amounts to assuming that, conditional on our historical and contemporary vectors of covariates, there are no unobserved confounders of the mediator-outcome relationship. We see this as a strong assumption, especially given that contemporary educational attainment can itself be considered a product of long-run socio-economic and political trajectories that we are unable to observe fully. In addition, any of these processes could in principle also leave a mark on today's political behaviour. The mediation results therefore serve the purpose of a decomposition exercise under this (strong) assumption rather than as a definitive account of how Protestant influence travels causally from the 19\textsuperscript{th} century to the present. We therefore only present evidence for our two mechanisms in \autoref{sec:mediation} of the Supplementary Material.

A number of additional analyses are reported in the Supplementary Material, including a sensitivity analysis based on \textcite{cinelli2020making}, and continuous-treatment entropy balancing.

\section{Results}

Our theory and argument lead to two empirical predictions. First, communes with larger historical Protestant populations should display lower contemporary support for the Radical Right. This negative association should hold even after accounting for the geographical, structural, and historical factors that shaped Protestant settlement itself. Crucially, it should not be entirely absorbed by contemporary educational attainment, since our theory posits a second, non-educational channel rooted in the political legacies of minority status -- a liberal Protestant ethic. Second, communes with larger historical Protestant populations should exhibit higher levels of human capital today, consistent with the long-run persistence of educational advantages associated with reformed religious practice. 

\subsection{Protestant Settlement and the Radical Right Vote}
Table \ref{tab:results_basic_ols} reports the baseline OLS estimates of the relationship between the 1839 commune-level share of Protestants and the first-round vote share of the \textit{Front National}/\textit{Rassemblement National} in the seven legislative elections held between 1993 and 2022. We progressively saturate the specification. We gradually add department fixed effects (1), historical controls only (2), contemporary controls only (3), both historical and contemporary covariates (4), and then, in turn, the remaining two columns (5 \& 6) add the share of the population holding a higher education degree and the share holding a \textit{baccalauréat}. Table \ref{tab:results_basic_ols_full} in the Supplementary Material includes all covariates.The results using sequential g-estimation, also gradually adding intermediate confounders, are shown in Table \ref{tab:results_seqg}. The coefficients on the 1839 Protestant share of the population are shown on Figure \ref{fig:coef_plot}.

\input{tables/baseline_ols_regression}

\input{tables/sequential_g_regression}

\setstretch{1.3}

\begin{figure}
    \centering
    \includegraphics[width=\textwidth]{coef_plot.png}
    \caption{Coefficient of 1839 Protestant share on RN vote share (Pooled)}
    \label{fig:coef_plot}
\end{figure}

The overall pattern is striking in its consistency. Across the ten model specifications, the estimated coefficient on the 1839 Protestant share moves only from $-0.027$ to $-0.024$, and remains statistically significant throughout. This stability is informative in itself. If the relationship were primarily driven by structural features of the regions where Protestantism took root (the Mediterranean mode of production, geographical isolation, the historical reach of the absolutist state etc.), we would expect the coefficient to attenuate or become insignificant once these features were accounted for. And at the same time, if it were primarily driven by contemporary socioeconomic differences between Protestant and non-Protestant communes, the contemporary controls (Model 3) should similarly absorb most of the effect either. Yet Table \ref{tab:results_basic_ols} yields that both of these scenarios do not turn out to be valid. We also see that post-treatment bias in the OLS estimates is likely to be slight.

The comparison between the full models before and after accounting for education (Model 4 compared to Models 5 and 6 in Table \ref{tab:results_basic_ols}, and Model 3 compared to Model 4 in Table \ref{tab:results_seqg}) is of special importance for our two-fold mechanism between human capital and minority status. Adding the contemporary share of higher education holders in the OLS specification produces, as expected, a strong negative association between education and Radical Right voting. We consider the coefficient of $-0.072$ in Model 5 and 6 as consistent with the broader literature on the educational divide in contemporary European electorates \parencite[see for example][]{ivarsflaten2012voting, abou-chadiBrahminLeftMerchant2021, hooghe2025field, hooghe2026does}. Importantly, the coefficient of the 1839 Protestant share is unchanged, both when controlling directly for education levels and when taking them into account in the sequential g-estimation specification. This is the empirical finding that our theory predicts: education carries some of the theorised relationship but adding the share of people with higher education or additionally high-school diploma holders again leaves the Protestant share estimate unchanged and statistically significant. A comparison of the two different education measures suggests that it is higher education specifically, rather than secondary schooling, capturing the educational pathway.

Figure \ref{fig:rn_over_time} reports the coefficient on the 1839 Protestant share from separate cross-sectional regressions estimated year-by-year, again varying the use of historical controls and contemporary education controls in both a naïve OLS setup and using sequential g-estimation. Over time, the negative relationship between historical share of Protestants and Radical Right voting is consistent. The estimated coefficient is negative and statistically significant in every single election year, and under every specification. Across the 63 model-year combinations, all 63 point estimates are negative, ranging from $-0.010$ to $-0.04$. %, with $p$-values that never exceed $0.006$ and that fall below $0.001$ in most cases.
%\footnote{The full year-by-year models can be found in Table \ref{tab:rn_by_year} of the Supplementary Material.} 

\begin{figure}[!htbp]
\centering
\includegraphics[width=\textwidth]{coef_per_year_seqg.png}
\caption{\Centering Coefficient of the Share of Protestants in 1839 on the Radical Right Vote Share, by Election Year. %The regression tables of the year-by-year models can be found in Table \ref{tab:rn_by_year} of the Supplementary Material.
}
\label{fig:rn_over_time}
\end{figure}

These findings suggest that the association of historical Protestant settlement and resistance to the Radical Right survived the rise of the FN under Jean-Marie Le Pen in the 1990s, the period of relative electoral marginalisation after the thwarted run for the presidency in 2002, the repositioning and attempts at ``dé-diabolisation" from 2011 onwards, and the renewed surge of the RN in the 2017 and 2022 election cycles which were marked by a rapidly increasing normalization of the Radical Right. The historical legacy reaching back to the early modern Wars of Religion appears to leave its association on French electoral geography.

While we do not want to over-interpret the oscillating pattern, since the year-by-year estimates are themselves noisy and the differences between adjacent years are often within each other's confidence intervals, the troughs in 2017 and 2022 deserve some attention. The coefficient is most negative in the years in which the Radical Right has been at its electorally strongest -- in 2017 when the FN reached the second round of the presidential election with Marine Le Pen, and in 2022, when the RN secured 89 seats in the National Assembly (81 more than in the previous term). This is what we would expect if the relevant mechanism is not a generic ``Protestant preference'' for non-conservative parties but something more specific to the politics of exclusionary nationalism: the historical commitments of formerly Protestant communities should become \textit{more} consequential, not less, as the Radical Right becomes more central to political competition and as the political stakes for minorities sharpen. Recent work on the activation of historical legacies has emphasised exactly this dynamic, with dormant patterns of political behaviour resurging in response to a heightened salience of the issues they originally addressed \parencite{wittenbergCruciblesPoliticalLoyalty2006, dinasInterruptedContinuitiesLocal2024, ochsnerActivatedHistoryCase2024, foukaCollectiveRemembrancePrivate2023}.

\subsection{Protestant Settlement and Contemporary Human Capital}

Tables \ref{tab:education_ols} and \ref{tab:education_seqg} examine the first link in the educational mechanism. In other words, they seek to assess whether historical Protestant settlement is associated with greater contemporary educational attainment. Here, the outcome is the commune-level share of the resident population holding a higher education degree. The specifications in Table \ref{tab:education_ols} mirror those of Models 2, 3, and 4 of Table \ref{tab:results_basic_ols}, whereas those in Table \ref{tab:education_seqg} correspond to Models 1, 2, and 3 of Table \ref{tab:results_seqg}.

The 1839 Protestant share is positively associated with contemporary higher education across all specifications, although the coefficient is more sensitive to the inclusion of controls than was the case for Radical Right voting in the subsection above - especially in the OLS setup. In the naïve OLS specification, and with historical controls alone, the estimate is $0.013$ ($p < 0.01$), whereas with contemporary controls alone, it drops to $0.009$ ($p < 0.05$). Adding both vectors of control, contemporary and historical, the coefficient sits at $0.008$ but is only significant at $p < 0.1$. This attenuation suggests that some of the simple correlation between the Protestant share and contemporary education reflects shared variation with present-day income, occupational composition, property values, and/or population structure. We regard this as unsurprising, since (past) education is highly endogenous to a wide set of contemporary characteristics that themselves correlate with both historical Protestantism and current educational outcomes. Indeed, when we take these intermediate confounders into account through sequential g-estimation, we find a more stable and significant coefficient of between $0.011$ and $0.013$. Overall, the estimate is stable in sign and order of magnitude across specifications. 

Taken together, these results suggest a modest but fragile contemporary educational association rather than a strong, broadly distributed legacy. If historical Protestantism only weakly predicts contemporary educational attainment, then the highly robust relationship between Protestant settlement and Radical Right resistance documented in Table \ref{tab:results_basic_ols} cannot principally be running through the human capital channel. The mediation analysis, reported in Section \ref{app:mediation} of the Supplementary Material and discussed in the robustness checks below, formalises this intuition. Contemporary education accounts for only a small fraction of the total association, with the bulk of the effect operating through other channels. We interpret this direct statistical effect, in line with our theoretical framework, as the political legacy of religious minority status.

\input{tables/protestants_higher_education}
\setstretch{1.3}


\subsection{Robustness Checks}
In Sections \ref{app:ols_regressions}, \ref{app:mediation}, \ref{app:sensitivity}, and \ref{app:weighting} of the Supplementary Material, we conduct a range of robustness checks. Figure \ref{fig:conley} shows the coefficients on the 1839 share of the Protestant population and corresponding confidence intervals when predicting contemporary support for the Radical Right in an OLS specification similar to Models 1 to 5 of Table \ref{tab:results_basic_ols_full} but using Conley standard errors at a variety of possible cutoffs, in order to take possible spatial correlation into account. We see that, for most model-cutoff pairs, the estimate remains significant (except for some models controlling exclusively for contemporary covariates).

We then estimate the models only with the communes explicitly mentioned in the original survey (Tables \ref{tab:results_original} and \ref{tab:educ_original}) and with the natural logarithm of the share of Protestants plus 1 (Tables \ref{tab:results_log} and \ref{tab:educ_log}). Across these specifications, the association between Protestant settlement and RN support remains negative and significant, while that between Protestant settlement and the contemporary level of tertiary education is still positive but not always significant at conventional levels. This is reinforced by the mediation analysis, which is intended as a way of comparing the relative role of two mechanisms, rather than to make a strong causal claim about indirect effects. The results shown in Table \ref{tab:mediation_results} suggest that contemporary human capital explains only a small fraction of the association between historical Protestant settlement and RN support.

To test the possibility that it might be religious competition with the Catholic Church, and not Protestantism itself, driving higher levels of human capital in areas where the Protestant population was greater, we interact the share of Protestants in 1839 with the share of constitutional priests in 1791, as a measure of (lack of) Catholic religiosity. Contrary to the expectations of this alternative account, we see in models 4 and 5 of Table \ref{tab:educ_full} of the Supplementary Material that the relationship between Protestantism and the share of the population with higher education today becomes stronger with the share of constitutional priests. It is in areas where most priests swore the Civil Constitution of the clergy (areas of weaker religiosity) that the association between Protestant settlement and human capital today is strongest. Conversely, in areas where most priests refused to do so in 1791 and where Catholicism was strongest, an increase in the Protestant share of the population is actually associated with a smaller share of the population with higher education.

As a further robustness check, we examined whether the negative association between historical Protestant settlement and Radical Right voting reflects a generic Protestant aversion to right-leaning parties or something more specific to the politics of the contemporary Radical Right. Table \ref{tab:mainstream_right_full} in the Supplementary Material re-estimates our six baseline specifications using the first-round vote share of the mainstream Right as the outcome variable. The relationship between the 1839 Protestant share and mainstream Right voting is also negative across all six specifications, but it is weaker in magnitude than the corresponding RN coefficients. In the fully saturated model with both historical and contemporary controls plus higher education and \textit{baccalauréat}, the estimated coefficient on the 1839 Protestant share is $-0.016$ for the mainstream Right compared to $-0.024$ for the Radical Right. The substantive pattern is therefore one in which historical Protestant settlement is associated with reduced support for the political right more broadly, but the relationship is most pronounced for the Radical Right. This is consistent with the broader French literature on the voting behaviour on the right and religiosity, which has documented a positive pull of Catholic religiosity toward the mainstream Right alongside no comparable, or even slightly negative, association with the Radical Right \parencite{lewis-beck2004, boyLelecteurSesRaisons1997a, mayerFrontNationalVote1998, michelatSystemeSymboliqueCatholique2015, overosDieuDansCommune2023}.

While we do not claim a causal effect here, the sensitivity analysis and the continuous-treatment entropy balancing shown in Sections \ref{app:sensitivity} and \ref{app:weighting} of the Supplementary Material provide a stronger basis for interpreting the pattern as substantive rather than purely compositional. We weight communes so that low- and high-Protestant areas are comparable on observed historical, geographic, and, in the extended specification, contemporary covariates; on the weighted sample, the negative association between historical Protestant presence and RN support remains clear. Sensitivity analysis further indicates that an omitted confounder would have to be extremely strongly related to both Protestant share and RN vote share (roughly 20 times the explanatory power of the benchmark covariate, mean income) to reduce the estimate to zero. Though our design is not causally identified, the observed association is not easily explained away by the most likely alternative accounts or by modest unobserved confounding.


\section{Discussion and Conclusion}
This paper started from a puzzle in the comparative scholarship on the political legacies of the Reformation. Protestantism has been linked to higher levels of human capital, to a distinctive work ethic, and, in the German case, to greater support for authoritarian nationalism in the early twentieth century. The French case does not fit this picture. French Protestants formed a small religious minority that lived through centuries of state-sponsored persecution, which led contemporary observers from \textcite{siegfriedGeographieElectoraleLArdeche1949} onwards to note their allegiance to the political left, as well as with Republican commitments to pluralism and freedom of religion. In this paper, we sought to give a systematic answer on how the confession most often associated in the social sciences with individualism, rigorous labour, and support for National Socialism came to be linked with resistance to exclusionary nationalism.

We argue that the political consequences of the Reformation cannot be reduced to a single mechanism. Two pathways connect historical Protestant settlement to contemporary political behaviour in France. The first runs through human capital, building on the long tradition stemming from \textcite{beckerWasWeberWrong2009} that links Protestantism to durable advantages in literacy and education. In turn, these factors are themselves negatively associated with contemporary support for Radical Right parties. The second runs through minority status. The historical experience of being a vulnerable religious minority in a Catholic political order generated political orientations centred on the defence of pluralism, religious freedom, and constitutional liberalism, in line with Rovny's \parencite*{rovnyEthnicMinoritiesPolitical2024, rovnyAntidoteBackslidingEthnic2023} account of minority groups as circumstantial liberals and with Braun's \parencite*{braunProtectorsPluralismReligious2019, braunReligiousMinoritiesResistance2016} evidence on inter-minority solidarity. These orientations should depress support for parties claiming to speak for a culturally homogeneous national majority, regardless of the educational composition of the communes in question.

Our empirical results support both predictions, though with slightly different degrees of confidence. We built a novel dataset on the historical distribution of French Protestants at the commune level, derived from a survey of the non-Catholic recognised churches conducted by the state in 1839. Combined with historical and contemporary commune-level data, this new data allows us to estimate the long-run association between Protestant settlement and contemporary political outcomes across the seven legislative elections held between 1993 and 2022. The relationship between the 1839 Protestant share and Radical Right voting is robust to extensive controls, stable across more than three decades of French politics, and largely orthogonal to contemporary educational composition. The association with contemporary higher education is positive but weaker, more sensitive to model specification, and less statistically reliable. A subsequent mediation analysis attributes only a small share of the total effect to the educational channel. We therefore conclude that the substantive picture is one in which the political legacy of French Protestantism reaches contemporary electoral behaviour primarily through channels other than human capital.

% discuss limitations

Several limitations of our analysis deserve emphasis. Our research design is observational and cross-sectional, and even with extensive historical and contemporary controls, we cannot fully rule out unobserved confounders that simultaneously shaped Protestant settlement centuries ago and continue to influence political behaviour today. Sensitivity analysis provides some reassurance, since a confounder strong enough to overturn our estimates would have to be roughly twenty times more explanatory than mean income, but it does not establish causal identification. We have also been deliberately cautious in the interpretation of the mediation analysis, which assumes sequential ignorability and rests on the additional assumption that contemporary educational attainment is, in the relevant sense, a clean mediator rather than a marker of further unobserved processes.

The empirical scope of the analysis is (for now) also incomplete. The data digitised so far covers only the Reformed (Calvinist) churches, leaving the Lutheran population concentrated in Alsace and the Montbéliard region outside our current sample. Incorporating the corresponding denomination is a priority for subsequent work. We also use commune-level electoral aggregates throughout, which means our findings speak to the geography of political behaviour rather than to individual-level mechanisms. Individual-level survey data, ideally linked to local Protestant presence, would allow a sharper test of whether the political orientations we identify operate at the level of individual descendants of Protestant communities, of the broader local political culture, or of some combination of the two.

A related limitation is the absence of current-day data on the spread of Protestantism in France. Without such contemporary measures, we are unable to empirically differentiate between two plausible transmission stories -- one in which the descendants of Protestant communities have themselves maintained distinctive political orientations across generations, and one in which historically Protestant areas developed durable local political cultures that now shape the behaviour of all residents regardless of individual religious lineage.

Despite these limitations, this paper contributes to the literature on the political legacies of religion in three substantive ways. First, they push back against the dominant framing of Protestantism in the social sciences as a confession whose political consequences are channelled primarily through human capital and a Weberian work ethic. The French case shows that Protestantism's political imprint can be substantial without being mediated by contemporary educational attainment. The dominant framings have largely emerged from cases in which Protestants formed the local majority and where the Reformation was a matter of state policy rather than dissident survival. In settings where Protestants instead lived as a persecuted minority, the political legacy may run through quite different channels, and existing frameworks risk obscuring rather than illuminating those legacies.

Second, our findings speak to a broader literature on the political consequences of historical minority status. Recent work has emphasised that minority groups can become durable defenders of liberal institutions, not because of doctrinal commitments but because such institutions constrain majoritarian power they cannot expect to wield themselves \parencite{rovnyAntidoteBackslidingEthnic2023, rovnyEthnicMinoritiesPolitical2024, braunProtectorsPluralismReligious2019}. Our results bring this argument to a setting in which the original minority experience is centuries old, the original persecution has long since ended, and the descendants of the minority population are now largely indistinguishable from the broader population on most observable dimensions. That historical minority status can still leave a measurable imprint on commune-level voting behaviour in twenty-first-century France suggests that minority-derived political orientations can outlive their original generating conditions by extraordinary periods of time.

Third, our results speak to ongoing debates on the rise of the Radical Right. The literature discusses general structural drivers of Radical Right support and emphasises factors such as economic decline, immigration, and cultural backlash. Similar to existing studies, our findings complement this work by highlighting the importance of deep historical local context. Communes that may look similar on present-day economic and demographic indicators can nevertheless display systematically different political responses to the contemporary Radical Right associated with legacies reaching back to Reformation. Such historical legacies need not fade as their original contexts recede but may instead be reactivated as their substantive concerns return to the centre of political contestation \parencite{wittenbergCruciblesPoliticalLoyalty2006, foukaCollectiveRemembrancePrivate2023, ochsnerActivatedHistoryCase2024, dinasInterruptedContinuitiesLocal2024}.

The political consequences of the Reformation in France are neither well captured by the standard Weberian and Beckerian frameworks nor reducible to a single mechanism. The educational legacy emphasised in much of the existing literature is real but appears to have substantially eroded under nearly two centuries, and it accounts for only a small share of the political distinctiveness of historically Protestant communes today. What remains is a striking and durable association between historical Protestant settlement and resistance to the Radical Right, which we read as the persistent political legacy of a minority experience that French Protestants themselves have long understood as central to their place in the Republic.

\clearpage

\printbibliography

\newpage

\input{appendix}



\end{document}
